\IfFileExists{stacks-project.cls}{%
\documentclass{stacks-project}
}{%
\documentclass{amsart}
}

% For dealing with references we use the comment environment
\usepackage{verbatim}
\newenvironment{reference}{\comment}{\endcomment}
%\newenvironment{reference}{}{}
\newenvironment{slogan}{\comment}{\endcomment}
\newenvironment{history}{\comment}{\endcomment}

% For commutative diagrams we use Xy-pic
\usepackage[all]{xy}

% We use 2cell for 2-commutative diagrams.
\xyoption{2cell}
\UseAllTwocells

% We use multicol for the list of chapters between chapters
\usepackage{multicol}

% This is generally recommended for better output
\usepackage{lmodern}
\usepackage[T1]{fontenc}

% For cross-file-references
\usepackage{xr-hyper}

% Package for hypertext links:
\usepackage{hyperref}

% For any local file, say "hello.tex" you want to link to please
% use \externaldocument[hello-]{hello}
\externaldocument[introduction-]{introduction}
\externaldocument[conventions-]{conventions}
\externaldocument[sets-]{sets}
\externaldocument[categories-]{categories}
\externaldocument[topology-]{topology}
\externaldocument[sheaves-]{sheaves}
\externaldocument[sites-]{sites}
\externaldocument[stacks-]{stacks}
\externaldocument[fields-]{fields}
\externaldocument[algebra-]{algebra}
\externaldocument[brauer-]{brauer}
\externaldocument[homology-]{homology}
\externaldocument[derived-]{derived}
\externaldocument[simplicial-]{simplicial}
\externaldocument[more-algebra-]{more-algebra}
\externaldocument[smoothing-]{smoothing}
\externaldocument[modules-]{modules}
\externaldocument[sites-modules-]{sites-modules}
\externaldocument[injectives-]{injectives}
\externaldocument[cohomology-]{cohomology}
\externaldocument[sites-cohomology-]{sites-cohomology}
\externaldocument[dga-]{dga}
\externaldocument[dpa-]{dpa}
\externaldocument[sdga-]{sdga}
\externaldocument[hypercovering-]{hypercovering}
\externaldocument[schemes-]{schemes}
\externaldocument[constructions-]{constructions}
\externaldocument[properties-]{properties}
\externaldocument[morphisms-]{morphisms}
\externaldocument[coherent-]{coherent}
\externaldocument[divisors-]{divisors}
\externaldocument[limits-]{limits}
\externaldocument[varieties-]{varieties}
\externaldocument[topologies-]{topologies}
\externaldocument[descent-]{descent}
\externaldocument[perfect-]{perfect}
\externaldocument[more-morphisms-]{more-morphisms}
\externaldocument[flat-]{flat}
\externaldocument[groupoids-]{groupoids}
\externaldocument[more-groupoids-]{more-groupoids}
\externaldocument[etale-]{etale}
\externaldocument[chow-]{chow}
\externaldocument[intersection-]{intersection}
\externaldocument[pic-]{pic}
\externaldocument[weil-]{weil}
\externaldocument[adequate-]{adequate}
\externaldocument[dualizing-]{dualizing}
\externaldocument[duality-]{duality}
\externaldocument[discriminant-]{discriminant}
\externaldocument[derham-]{derham}
\externaldocument[local-cohomology-]{local-cohomology}
\externaldocument[algebraization-]{algebraization}
\externaldocument[curves-]{curves}
\externaldocument[resolve-]{resolve}
\externaldocument[models-]{models}
\externaldocument[functors-]{functors}
\externaldocument[equiv-]{equiv}
\externaldocument[pione-]{pione}
\externaldocument[etale-cohomology-]{etale-cohomology}
\externaldocument[proetale-]{proetale}
\externaldocument[relative-cycles-]{relative-cycles}
\externaldocument[more-etale-]{more-etale}
\externaldocument[trace-]{trace}
\externaldocument[crystalline-]{crystalline}
\externaldocument[spaces-]{spaces}
\externaldocument[spaces-properties-]{spaces-properties}
\externaldocument[spaces-morphisms-]{spaces-morphisms}
\externaldocument[decent-spaces-]{decent-spaces}
\externaldocument[spaces-cohomology-]{spaces-cohomology}
\externaldocument[spaces-limits-]{spaces-limits}
\externaldocument[spaces-divisors-]{spaces-divisors}
\externaldocument[spaces-over-fields-]{spaces-over-fields}
\externaldocument[spaces-topologies-]{spaces-topologies}
\externaldocument[spaces-descent-]{spaces-descent}
\externaldocument[spaces-perfect-]{spaces-perfect}
\externaldocument[spaces-more-morphisms-]{spaces-more-morphisms}
\externaldocument[spaces-flat-]{spaces-flat}
\externaldocument[spaces-groupoids-]{spaces-groupoids}
\externaldocument[spaces-more-groupoids-]{spaces-more-groupoids}
\externaldocument[bootstrap-]{bootstrap}
\externaldocument[spaces-pushouts-]{spaces-pushouts}
\externaldocument[spaces-chow-]{spaces-chow}
\externaldocument[groupoids-quotients-]{groupoids-quotients}
\externaldocument[spaces-more-cohomology-]{spaces-more-cohomology}
\externaldocument[spaces-simplicial-]{spaces-simplicial}
\externaldocument[spaces-duality-]{spaces-duality}
\externaldocument[formal-spaces-]{formal-spaces}
\externaldocument[restricted-]{restricted}
\externaldocument[spaces-resolve-]{spaces-resolve}
\externaldocument[formal-defos-]{formal-defos}
\externaldocument[defos-]{defos}
\externaldocument[cotangent-]{cotangent}
\externaldocument[examples-defos-]{examples-defos}
\externaldocument[algebraic-]{algebraic}
\externaldocument[examples-stacks-]{examples-stacks}
\externaldocument[stacks-sheaves-]{stacks-sheaves}
\externaldocument[criteria-]{criteria}
\externaldocument[artin-]{artin}
\externaldocument[quot-]{quot}
\externaldocument[stacks-properties-]{stacks-properties}
\externaldocument[stacks-morphisms-]{stacks-morphisms}
\externaldocument[stacks-limits-]{stacks-limits}
\externaldocument[stacks-cohomology-]{stacks-cohomology}
\externaldocument[stacks-perfect-]{stacks-perfect}
\externaldocument[stacks-introduction-]{stacks-introduction}
\externaldocument[stacks-more-morphisms-]{stacks-more-morphisms}
\externaldocument[stacks-geometry-]{stacks-geometry}
\externaldocument[moduli-]{moduli}
\externaldocument[moduli-curves-]{moduli-curves}
\externaldocument[examples-]{examples}
\externaldocument[exercises-]{exercises}
\externaldocument[guide-]{guide}
\externaldocument[desirables-]{desirables}
\externaldocument[coding-]{coding}
\externaldocument[obsolete-]{obsolete}
\externaldocument[fdl-]{fdl}
\externaldocument[index-]{index}

% Theorem environments.
%
\theoremstyle{plain}
\newtheorem{theorem}[subsection]{Theorem}
\newtheorem{proposition}[subsection]{Proposition}
\newtheorem{lemma}[subsection]{Lemma}

\theoremstyle{definition}
\newtheorem{definition}[subsection]{Definition}
\newtheorem{example}[subsection]{Example}
\newtheorem{exercise}[subsection]{Exercise}
\newtheorem{situation}[subsection]{Situation}

\theoremstyle{remark}
\newtheorem{remark}[subsection]{Remark}
\newtheorem{remarks}[subsection]{Remarks}

\numberwithin{equation}{subsection}

% Macros
%
\def\lim{\mathop{\mathrm{lim}}\nolimits}
\def\colim{\mathop{\mathrm{colim}}\nolimits}
\def\Spec{\mathop{\mathrm{Spec}}}
\def\Hom{\mathop{\mathrm{Hom}}\nolimits}
\def\Ext{\mathop{\mathrm{Ext}}\nolimits}
\def\SheafHom{\mathop{\mathcal{H}\!\mathit{om}}\nolimits}
\def\SheafExt{\mathop{\mathcal{E}\!\mathit{xt}}\nolimits}
\def\Sch{\mathit{Sch}}
\def\Mor{\mathop{\mathrm{Mor}}\nolimits}
\def\Ob{\mathop{\mathrm{Ob}}\nolimits}
\def\Sh{\mathop{\mathit{Sh}}\nolimits}
\def\NL{\mathop{N\!L}\nolimits}
\def\CH{\mathop{\mathrm{CH}}\nolimits}
\def\proetale{{pro\text{-}\acute{e}tale}}
\def\etale{{\acute{e}tale}}
\def\QCoh{\mathit{QCoh}}
\def\Ker{\mathop{\mathrm{Ker}}}
\def\Im{\mathop{\mathrm{Im}}}
\def\Coker{\mathop{\mathrm{Coker}}}
\def\Coim{\mathop{\mathrm{Coim}}}

% Boxtimes
%
\DeclareMathSymbol{\boxtimes}{\mathbin}{AMSa}{"02}

%
% Macros for moduli stacks/spaces
%
\def\QCohstack{\mathcal{QC}\!\mathit{oh}}
\def\Cohstack{\mathcal{C}\!\mathit{oh}}
\def\Spacesstack{\mathcal{S}\!\mathit{paces}}
\def\Quotfunctor{\mathrm{Quot}}
\def\Hilbfunctor{\mathrm{Hilb}}
\def\Curvesstack{\mathcal{C}\!\mathit{urves}}
\def\Polarizedstack{\mathcal{P}\!\mathit{olarized}}
\def\Complexesstack{\mathcal{C}\!\mathit{omplexes}}
% \Pic is the operator that assigns to X its picard group, usage \Pic(X)
% \Picardstack_{X/B} denotes the Picard stack of X over B
% \Picardfunctor_{X/B} denotes the Picard functor of X over B
\def\Pic{\mathop{\mathrm{Pic}}\nolimits}
\def\Picardstack{\mathcal{P}\!\mathit{ic}}
\def\Picardfunctor{\mathrm{Pic}}
\def\Deformationcategory{\mathcal{D}\!\mathit{ef}}
\usepackage{cleveref}
\crefname{originallemma}{original lemma}{original lemmas}
\Crefname{originallemma}{Original lemma}{Original lemmas}
\makeatletter
\newlabel{lemma-groupoid-on-field-explain-points}{{lemma-groupoid-on-field-explain-points}{source}{}{sourcepointer-1}{}}
\newlabel{lemma-groupoid-on-field-explain-points@cref}{{[originallemma][0][]lemma-groupoid-on-field-explain-points}{[1][0][]source}{}{}{}}
\makeatother
\crefname{originallemma}{original lemma}{original lemmas}
\Crefname{originallemma}{Original lemma}{Original lemmas}
\makeatletter
\newlabel{lemma-groupoid-on-field-move-point}{{lemma-groupoid-on-field-move-point}{source}{}{sourcepointer-2}{}}
\newlabel{lemma-groupoid-on-field-move-point@cref}{{[originallemma][0][]lemma-groupoid-on-field-move-point}{[1][0][]source}{}{}{}}
\makeatother
\crefname{originallemma}{original lemma}{original lemmas}
\Crefname{originallemma}{Original lemma}{Original lemmas}
\makeatletter
\newlabel{lemma-restrict-groupoid-on-field}{{lemma-restrict-groupoid-on-field}{source}{}{sourcepointer-3}{}}
\newlabel{lemma-restrict-groupoid-on-field@cref}{{[originallemma][0][]lemma-restrict-groupoid-on-field}{[1][0][]source}{}{}{}}
\makeatother
\crefname{originallemma}{original lemma}{original lemmas}
\Crefname{originallemma}{Original lemma}{Original lemmas}
\makeatletter
\newlabel{varieties-lemma-bijection-irreducible-components}{{038F}{source}{}{sourcepointer-4}{}}
\newlabel{varieties-lemma-bijection-irreducible-components@cref}{{[originallemma][0][]038F}{[1][0][]source}{}{}{}}
\makeatother
\crefname{originallemma}{original lemma}{original lemmas}
\Crefname{originallemma}{Original lemma}{Original lemmas}
\makeatletter
\newlabel{varieties-lemma-separably-closed-irreducible}{{020J}{source}{}{sourcepointer-5}{}}
\newlabel{varieties-lemma-separably-closed-irreducible@cref}{{[originallemma][0][]020J}{[1][0][]source}{}{}{}}
\makeatother
\crefname{originallemma}{original lemma}{original lemmas}
\Crefname{originallemma}{Original lemma}{Original lemmas}
\makeatletter
\newlabel{varieties-lemma-orbit-irreducible-components}{{04KY}{source}{}{sourcepointer-6}{}}
\newlabel{varieties-lemma-orbit-irreducible-components@cref}{{[originallemma][0][]04KY}{[1][0][]source}{}{}{}}
\makeatother
\title{Unique irreducible component through the identity of a groupoid over a field, geometrically irreducible}
\author{The Stacks Project Authors}
\date{stacks-project:more-groupoids.tex:local-snapshot-20261012; GNU FDL 1.2 or later}
\begin{document}
\maketitle
\setcounter{subsection}{1}
\noindent\textit{Editorial note.} All delivered mathematical text is the author's own, copied from the pinned Stacks Project chapter file more-groupoids.tex; no statement or proof step was rewritten, reordered or removed. The counted span is the author environment and its own proof environment. References inside the span are handed to the assembler as navigation stubs that carry the canonical tag of the cited tag only; whether each of them is load-bearing is recorded on the stub itself.\par\medskip
\noindent\textit{Allowed standard prerequisites (editorial).}\begin{itemize}
\item Base change making the source and target maps isomorphisms at a chosen point of the restricted groupoid (more-groupoids chapter, Lemma lemma-groupoid-on-field-explain-points, no canonical tag in the snapshot; no-tag pointer to the source label line; the item is also collected on this lane as L22505). cited at proof lines 1459-1460. Provides the field extension and point r' after which s, t are isomorphisms, enabling the move-point reduction.
\item The restriction of a groupoid on a field to a field extension is faithfully flat on local rings (more-groupoids chapter, Lemma lemma-restrict-groupoid-on-field, no canonical tag in the snapshot; no-tag pointer to the source label line). cited at proof lines 1461-1463 and 1494-1495. Faithfully flat local maps transfer minimal primes and irreducible components between R and its base changes, in both reduction directions.
\item Automorphism moving the identity to a point when source and target are isomorphisms (more-groupoids chapter, Lemma lemma-groupoid-on-field-move-point, no canonical tag in the snapshot; no-tag pointer to the source label line; the item is also collected on this lane as L22506). cited at proof lines 1466-1468. Reduces part (1) at an arbitrary point r to part (1) at the identity, hence to part (2).
\item Bijection between irreducible components and generic points used for products over a separably closed field (varieties chapter, Lemma varieties-lemma-bijection-irreducible-components, canonical tag 038F). cited at proof lines 1474-1476. Input to the irreducibility of the product of two components through e in the separably closed case.
\item Over a separably algebraically closed field the product of irreducible schemes with a common point is irreducible (varieties chapter, Lemma varieties-lemma-separably-closed-irreducible, canonical tag 020J). cited at proof lines 1474-1477. Gives irreducibility of the fibre product of two components through e.
\item Irreducible components of a base change over a field form a Galois orbit (varieties chapter, Lemma varieties-lemma-orbit-irreducible-components, canonical tag 04KY). cited at proof lines 1504-1506. Every component of the base change of Z is a Galois translate of Z'', forcing all of them equal to the unique component through e'' and giving geometric irreducibility.
\end{itemize}\medskip
\ifdefined\hypertarget\hypertarget{sourcepointer-1}{}\fi
\noindent\textit{Source reference.} Original lemma lemma-groupoid-on-field-explain-points: sources/stacks-more-groupoids/pinned-original-source.tex, source lines 1258--1258. Cited at proof lines 1459-1460: base change making s, t isomorphisms at the chosen point.\par\medskip
\ifdefined\hypertarget\hypertarget{sourcepointer-2}{}\fi
\noindent\textit{Source reference.} Original lemma lemma-groupoid-on-field-move-point: sources/stacks-more-groupoids/pinned-original-source.tex, source lines 1329--1329. Cited at proof lines 1466-1468: automorphism moving e to r reduces part (1) to part (2).\par\medskip
\ifdefined\hypertarget\hypertarget{sourcepointer-3}{}\fi
\noindent\textit{Source reference.} Original lemma lemma-restrict-groupoid-on-field: sources/stacks-more-groupoids/pinned-original-source.tex, source lines 1219--1219. Cited at proof lines 1461-1463 and 1494-1495: faithfully flat local ring maps transfer the statements between R and its base changes.\par\medskip
\ifdefined\hypertarget\hypertarget{sourcepointer-4}{}\fi
\noindent\textit{Source reference.} Original lemma 038F: sources/stacks-chow/tags.txt, source lines 4234--4234. Cited at proof lines 1474-1476: input to irreducibility of the product of two components in the separably closed case.\par\medskip
\ifdefined\hypertarget\hypertarget{sourcepointer-5}{}\fi
\noindent\textit{Source reference.} Original lemma 020J: sources/stacks-chow/tags.txt, source lines 2636--2636. Cited at proof lines 1474-1477: the fibre product of two components through e is irreducible.\par\medskip
\ifdefined\hypertarget\hypertarget{sourcepointer-6}{}\fi
\noindent\textit{Source reference.} Original lemma 04KY: sources/stacks-chow/tags.txt, source lines 5989--5989. Cited at proof lines 1504-1506: components of the base change of Z are Galois translates of Z'', giving geometric irreducibility.\par\medskip
\begin{lemma}
\label{lemma-groupoid-on-field-geometrically-irreducible}
Let $S$ be a scheme. Let $(U, R, s, t, c)$ be a groupoid scheme
over $S$. Assume $U = \Spec(k)$ with $k$ a field.
By abuse of notation denote $e \in R$ the image of the identity
morphism $e : U \to R$. Then
\begin{enumerate}
\item every local ring $\mathcal{O}_{R, r}$ of $R$ has a unique
minimal prime ideal,
\item there is exactly one irreducible component $Z$ of $R$
passing through $e$, and
\item $Z$ is geometrically irreducible over $k$ via either
$s$ or $t$.
\end{enumerate}
\end{lemma}

\begin{proof}
Let $r \in R$ be a point.
In this proof we will use the correspondence between irreducible components
of $R$ passing through a point $r$ and minimal primes of the local
ring $\mathcal{O}_{R, r}$ without further mention.
Choose $k \subset k'$ and $r' \in R'$ as in
Lemma \ref{lemma-groupoid-on-field-explain-points}.
Note that $\mathcal{O}_{R, r} \to \mathcal{O}_{R', r'}$
is faithfully flat and local, see
Lemma \ref{lemma-restrict-groupoid-on-field}.
Hence the result for $r' \in R'$ implies the result for $r \in R$.
In other words we may assume that $s, t : k \to \kappa(r)$
are isomorphisms. By
Lemma \ref{lemma-groupoid-on-field-move-point}
there exists an automorphism moving $e$ to $r$.
Hence we may assume $r = e$, i.e., part (1) follows from part (2).

\medskip\noindent
We first prove (2) in case $k$ is separably algebraically closed.
Namely, let $X, Y \subset R$ be irreducible components
passing through $e$. Then by
Varieties, Lemma \ref{varieties-lemma-bijection-irreducible-components} and
\ref{varieties-lemma-separably-closed-irreducible}
the scheme $X \times_{s, U, t} Y$ is irreducible as well.
Hence $c(X \times_{s, U, t} Y) \subset R$ is an irreducible subset.
We claim it contains both $X$ and $Y$ (as subsets of $R$).
Namely, let $T$ be the spectrum of a field. If $x : T \to X$ is a $T$-valued
point of $X$, then $c(x, e \circ s \circ x) = x$ and $e \circ s \circ x$
factors through $Y$ as $e \in Y$. Similarly for points of $Y$.
This clearly implies that $X = Y$, i.e., there is a unique irreducible
component of $R$ passing through $e$.

\medskip\noindent
Proof of (2) and (3) in general.
Let $k \subset k'$ be a separable algebraic closure, and
let $(U', R', s', t', c')$ be the restriction of
$(U, R, s, t, c)$ via $\Spec(k') \to \Spec(k)$.
By the previous paragraph there is exactly one irreducible
component $Z'$ of $R'$ passing through $e'$.
Denote $e'' \in R \times_{s, U} U'$ the base change of $e$.
As $R' \to R \times_{s, U} U'$ is faithfully flat, see
Lemma \ref{lemma-restrict-groupoid-on-field},
and $e' \mapsto e''$ we see that there is exactly one
irreducible component $Z''$ of $R \times_{s, k} k'$ passing
through $e''$. This implies, as $R \times_k k' \to R$ is faithfully
flat, that there is exactly one irreducible component $Z$ of $R$
passing through $e$. This proves (2).

\medskip\noindent
To prove (3) let $Z''' \subset R \times_k k'$ be an arbitrary
irreducible component of $Z \times_k k'$. By
Varieties, Lemma \ref{varieties-lemma-orbit-irreducible-components}
we see that $Z''' = \sigma(Z'')$ for some $\sigma \in \text{Gal}(k'/k)$.
Since $\sigma(e'') = e''$ we see that $e'' \in Z'''$ and hence
$Z''' = Z''$. This means that $Z$ is geometrically irreducible
over $\Spec(k)$ via the morphism $s$.
The same argument implies that $Z$ is geometrically irreducible
over $\Spec(k)$ via the morphism $t$.
\end{proof}

\end{document}
